\section{Improved Allocator Design}
\label{sec:improved-design}

\subsection{Design Goal}

The goal of v8.0.0 was not to completely replace the allocator policy introduced by v7.0.0, but to improve
the part of that design that too clumsy. In v7.0.0, the allocator already had segregated bins, thread-local
small-block caches, batch refill, deferred coalescing, and bounded retained cache memory. However, the problem
was that cache growth was still driven mainly by refill behavior. That policy could react to demand,
but it did so through an indirect signal.

The design goal in v8.0.0 was therefore to keep the same general architecture while changing how
the allocator decides which small-size classes deserve more cached capacity. Instead of treating repeated
refill events as the main indicator of importance, v8.0.0 measures live demand directly for exact reusable
slab sizes and uses that information to retune cache targets. This makes the allocator more workload-aware
without giving up the drop-in \texttt{new}/\texttt{delete} interface as they are part of the project's
deliverables. That direction is consistent with the broader motivation for custom allocators: once the
allocation pattern of a program is better understood, the allocator can be adapted to that pattern rather
than paying the full generality cost of a one-size-fits-all design~\cite{rakosCustomAllocators2021}.

\subsection{Baseline: What v7.0.0 Already Provided}

Version~7.0.0 inherited most of the important structural changes introduced between v3.0.0 and v6.0.0. Free
blocks were already grouped into segregated bins, which reduced the need to scan the entire heap for every
request. Small allocations could often be served from a thread-local cache, which reduced repeated global
locking on hot size classes. When those caches missed, the allocator could refill them in batches, gradually
decreasing the cost of central-pool access across several future allocations.

Version~7.0.0 no longer merged adjacent free blocks on every deallocation. Instead, it retained
the deferred-coalescing strategy from v5.0.0 and v6.0.0, so adjacent blocks were merged only when an allocation
request could not be satisfied directly. That design trades some fragmentation risk for lower steady-state
deallocation cost, a tradeoff that is common in dynamic allocation systems~\cite{wilson1995}. V7.0.0 also
limited retained cache bytes per thread, so cache growth was bounded even when a workload exercised the same
classes heavily.

These features matter because they define what V8.0.0 did \textit{not} need to change. The allocator already
had a practical free-structure layout, a central best-fit style search within bins, a cache for small hot
objects, and correctness machinery such as canaries and ownership checks. The main issue was not missing
functionality, but that one piece of policy, cache sizing, was still too blunt.

\subsection{Weakness in v7.0.0}

The weakness in v7.0.0 was that cache growth was based mostly on refill frequency. If a particular bin missed
repeatedly, the allocator interpreted that pattern as evidence that the class should keep more cached blocks.
That idea is sensible at a high level, but refill count is not the same thing as measured working-set demand.
A class can refill frequently because of a short burst, irregular reuse, or because the cached
capacity is briefly too small even though the long-term live set is modest.

This distinction matters because refill count is only an indirect signal. It tells the allocator that pressure
occurred, but not how many blocks were actually needed live at once. In practice, that means a hot phase can
leave a class overgrown after the workload shifts, while another class may deserve cached capacity for reasons
that are not well captured by miss events alone. The result overly reactive policy that might not be reacting
for the right reasons. V8.0.0 addresses exactly this problem by replacing refill-driven growth with measured
occupancy-driven retuning.

\subsection{High-Level V8.0.0 Strategy}

At a high level, V8.0.0 keeps the cache policy of V7.0.0. The allocator still uses segregated small and large
bins, still serves exact small slabs from a thread-local cache when possible, and still falls back to the
central pool on cache misses. But V8.0.0 adds explicit accounting for live exact-slab allocations and uses
that accounting to retune cache targets per size class.

The implementation also makes the allocator more observable. V8.0.0 introduces a metrics surface that records
cache hits, misses, searches, refills, flushes, tuning events, and mutex acquisitions. This does not directly
make the allocator faster, but it makes the allocator easier to analyze. That is a real design improvement in
the context of this assignment because elapsed time alone cannot explain why one workload improved while
another regressed.

\subsection{New Data Structures and Metrics}

The most visible interface addition in V8.0.0 is the \texttt{AllocatorMetrics} structure. It exposes counters
for thread-cache hits and misses, central-pool searches, searched free-list nodes, cache refills, cache
flushes, target-driven and byte-limit-driven flush events, coalescing on allocation failure, central mutex
acquisitions, and cache-policy tuning activity. Essentially turning the allocator into something that can
be measured rather than only timed.

The key internal addition is \texttt{CacheAccounting}. Each accounting record stores live-block counts for all
small-bin classes. When a thread cache is registered, it receives one of these accounting records and uses it
to track how many blocks of each exact slab size are live at the same time. Instead of inferring demand from
miss history, V8.0.0 can observe demand directly. The code also extends \texttt{ThreadCache} with per-bin arrays
for target sizes, peak live counts, and countdowns until the next tuning event.

Metrics are optionally enabled through a compile-time macro. That decision is important because profiling
counters are useful for analysis, but they can interfere with timing runs and skew data. The design therefore
separates observability from baseline execution cost.

\subsection{Exact Slab Classification}

Another important change in V8.0.0 is that the small-object cache is restricted to exact slab classes. The
helper \texttt{smallSlabSize} computes whether an allocation can be represented as one of these small exact
slab sizes after metadata, alignment, and quantum rounding are accounted for. If so, the allocator may serve it
from the thread cache; if not, the request follows the central-pool path.

This is more precise than treating every request below a broad threshold as interchangeable. Exact slab
classification means the cache is tuned around the actual reusable unit, not just the requested payload size.
That keeps the policy cleaner and reduces the chance that slightly different requests contaminate the same
reuse class. It also aligns naturally with the new accounting model, since occupancy is now tracked for the
slab sizes that the allocator truly reuses.

\subsection{Adaptive Cache Target Tuning}

Adaptive target tuning is the core improvement in V8.0.0. For each exact cached slab class, the allocator
tracks the current live count and the peak live count seen during a tuning window. After a fixed number of
allocations in that class, the allocator resets the cache target to a bounded function of observed demand:

\begin{equation}
\begin{aligned}
\text{new target} = {}& \operatorname{clamp}(\text{peak live blocks} + \text{cushion},\\
& \text{minimum},\ \text{maximum}).
\end{aligned}
\end{equation}

In the current implementation, the minimum cached target is eight blocks, the cushion is four blocks,
the retuning interval is 256 allocations in the class, and the maximum cached target is 512 blocks. This
policy is simple, but it is already more informative than refill-count doubling. A class grows because the
program actually kept that many blocks live, not only because a short miss burst happened to occur.

This retuning model still uses a heuristic, so it should not be described as optimal. However, it is a more
direct approximation to working-set size than the V7.0.0 rule. It also fits the broader principle that cache
effectiveness depends on matching capacity to the reuse pattern of the workload, not just reacting to misses
after the fact~\cite{hennessy2019}.

\subsection{Allocation and Deallocation Path Changes}

On the allocation path, V8.0.0 still begins by validating size and alignment and by registering the thread
cache if needed. If the request belongs to an exact small slab class, the allocator attempts a thread-cache
hit first. A successful hit no longer only returns a block; it also increments the accounting state through
\texttt{exactAllocationCounter}. If the request misses the thread cache, the allocator records the miss,
searches the central pool, and refills the thread cache when the request still belongs to an exact reusable
slab size.

The deallocation path was updated to mirror that model. When a block came from an exact cached slab class,
deallocation decrements the same live-count structure before the block is returned to the thread cache. This
is essential. Without deallocation-side accounting, the allocator could see when pressure rises, but not when
a class cools down. Updating both paths keeps the model synchronized with actual object lifetime.

Outside those changes, the general slow path remains familiar. Free-list insertion and removal still happen
in the central allocator, and coalescing still occurs only when allocation pressure requires it. That
continuity is useful because it isolates the effect of the new policy rather than mixing it with a wholesale
redesign of the free structure.

\subsection{What Stayed the Same}

Several important parts of the allocator did not change in v8.0.0. The allocator is still a singleton memory
pool behind global \texttt{new}, \texttt{new[]}, \texttt{delete}, and \texttt{delete[]}. It still uses
segregated small and large bins, retains a thread-local small-object cache, uses a central mutex for
global free-list work, and preserves canary-based corruption detection and ownership validation.

This continuity matters for interpretation. It means that v8.0.0 should not be described as an entirely new
allocator, and any measured behavior should not be attributed to a change in the whole allocation architecture.
Instead, the main change is a policy refinement inside an already established framework. That makes both the
design argument and the later experimental comparison cleaner.

\subsection{Why These Changes Were Chosen}

These changes were chosen for three practical reasons. First, they improve signal quality. Measured live
occupancy is closer to the true working set than refill count, so it is a better basis for deciding how much
cache capacity a class deserves. Second, they improve observability. By recording allocator events directly,
v8.0.0 makes it possible to explain results with evidence rather than only with conjecture. Third, they
preserve compatibility. The allocator remains reusable and fetchable by another project, and it still acts
as a drop-in replacement for the standard allocation interface.

There is also a pragmatic engineering reason for preferring this approach over a full rewrite. A focused
policy change is easier to validate, easier to compare against V7.0.0, and less likely to introduce unrelated
regressions. For this assignment, that is a better trade than introducing an entirely different
memory-management strategy and then trying to separate the effects afterward.

\subsection{Expected Effect on the UsingMatrixClass Workloads}

The expected effect of v8.0.0 depends on the workload. The matrix-array workload repeatedly creates and
replaces many objects with highly regular sizes. That regularity should give the exact-slab cache policy a
fair opportunity to discover stable hot classes and keep those classes warm. Uniform matrix nodes add more
lifetime churn, but still retain a relatively regular size pattern, so they should also benefit when live
occupancy remains concentrated in a narrow set of bins.

Nested matrix stress is a more mixed case. It uses several matrix sizes and repeatedly rebuilds subtrees, so
it exercises more classes and more transitions between them. This is a useful workload for testing whether
adaptive targets actually follow demand rather than simply inflating over time. The linked-list workload is
the most difficult case because it mixes many tiny nodes with additional dynamic arrays and highly active
updates. Even if v8.0.0 does not always improve this case, the policy is still motivated, since it should be
less likely than V7.0.0 to keep cold classes oversized after a burst.

The key point is that v8.0.0 was designed to improve policy quality, not to guarantee universal speedup. A
workload that benefits from stable exact slab reuse is a natural candidate for improvement. A workload
dominated by tiny-object metadata costs or central-pool behavior may still expose limitations.

\subsection{Validation and Test Support}

The repository includes tests that support both correctness and the new policy claims. The cache-focused
tests verify that reused small allocations can be served from the thread cache without forcing unnecessary
central-pool work. The cache-policy tests verify that the adaptive mechanism actually tunes targets and is
capable of both increasing and decreasing them over time. General memory-pool tests still cover best-fit reuse,
coalescing behavior, multithreaded use, and canary detection, while global allocation tests verify that the
allocator remains a correct replacement for ordinary \texttt{new} and \texttt{delete}.

This matters for the report because it shows that the improved design was not only described conceptually.
The design was implemented in a way that exposes measurable policy behavior and was tested at both the
allocator-core level and the integration level.

\subsection{Limitations}

V8.0.0 still has important limitations. The central free structure remains protected by one mutex, so global
allocator work is still serialized. Large allocations still depend on the same central allocator structure.
Metadata and canary checks remain present, which is useful for safety but still costs time and space. The
retuning policy itself is empirical and bounded by a few fixed constants, so it should be treated as a
practical heuristic rather than a theoretically optimal rule.

These limitations should be stated clearly because they define the scope of the contribution. V8.0.0
improves cache sizing and allocator observability. It does not eliminate all contention, all fragmentation
costs, or all workload-specific regressions.

\subsection{Summary}

Version~8.0.0 is best understood as a measured refinement of V7.0.0. It preserves the allocator family that
earlier versions built, but replaces refill-driven cache growth with occupancy-driven retuning for exact
reusable small slabs. At the same time, it introduces explicit metrics and accounting so allocator behavior
can be studied directly. The result is not a brand-new allocator architecture, but a more informed and more
analyzable one. That makes it a meaningful improvement for this project and a defensible basis for the
experimental comparisons in the next section.
